% GENERATED FILE. Edit canonical lessons, then run npm run generate:books.
% canonical-source-hash: fnv1a64:48717984b4c3da6d
% canonical-lessons: GU-C169-bhakhri, GU-C169-khicdi, GU-C169-papad, GU-C169-kadai, GU-C169-dhaknu

\chapter{Khichdi, bhakhri and papad}
\label{ch:gu-khichdi-bhakhri-and-papad}
\hlchaptermodality{169}

\begin{chapteropening}
\textbf{By the end of this chapter:} \emph{I can name everyday Gujarati dishes and the pans they are made in.}

Complete the last lesson of chapter 169: I can name everyday Gujarati dishes and the pans they are made in.
\end{chapteropening}

\section[\texorpdfstring{\gu{ભાખરી} (bhākhrī)}{ભાખરી (bhākhrī)}]{\gu{ભાખરી} (bhākhrī) — bhakhri, thick flatbread}

\label{lesson:GU-C169-bhakhri}



\begin{quote}
Before the new one: say the Gujarati for griddle, then the Gujarati for to roast.
\end{quote}

\subsection*{You'll want to know: \gu{ભાખરી}}
\textbf{\gu{ભાખરી}} — \emph{bhākhrī} — \textquotedblleft{}bhakhri, thick flatbread\textquotedblright{}.

\begin{quote}
\textbf{\gu{દાદા} \gu{સવારે} \gu{ભાખરી} \gu{ખાય} \gu{છે}.} — \emph{dādā savāre bhākhrī khāy che} — \textquotedblleft{}Grandfather eats bhakhri in the morning.\textquotedblright{}
\end{quote}

\begin{cousinweb}
\textbf{\gu{ભાખરી}} (\emph{bhākhrī}) is inherited. CDIAL lists the kindred masculine \emph{bhākhro}, \textquotedblleft{}lump of bread\textquotedblright{}, under \emph{*bhakkha}, \textquotedblleft{}lump\textquotedblright{}, a reconstructed form, though not \emph{bhākhrī} itself (CDIAL 9330).
\end{cousinweb}

\begin{grammarlens}[title={What you've built}]
One more word for everyday life at A2.
\end{grammarlens}

\subsection*{Your turn}
\begin{itemize}
\raggedright
  \item \emph{Say it:} \emph{bhākhrī}
  \item \emph{Say it:} \emph{bhākhrī}, once more
  \item \emph{Say it:} \emph{dādā savāre bhākhrī khāy che} — Grandfather eats bhakhri in the morning.
  \item \emph{Recall:} say the Gujarati for griddle, then the Gujarati for to roast, then say \emph{bhākhrī} again
\end{itemize}

\begin{tcolorbox}[breakable,colback=teal!4,colframe=teal!35!black,title={Before you move on}]
What does \gu{ભાખરી} mean? (\textquotedblleft{}Bhakhri, thick flatbread\textquotedblright{}.) Say it once more.
\end{tcolorbox}
\section[\texorpdfstring{\gu{ખીચડી} (khīcḍī)}{ખીચડી (khīcḍī)}]{\gu{ખીચડી} (khīcḍī) — khichdi, rice and lentils}

\label{lesson:GU-C169-khicdi}



\begin{quote}
Before the new one: say the Gujarati for to roast, then the Gujarati for bhakhri.
\end{quote}

\subsection*{You'll want to know: \gu{ખીચડી}}
\textbf{\gu{ખીચડી}} — \emph{khīcḍī} — \textquotedblleft{}khichdi, rice and lentils\textquotedblright{}.

\begin{quote}
\textbf{\gu{રાત્રે} \gu{અમે} \gu{ખીચડી} \gu{જમીએ} \gu{છીએ}.} — \emph{rātre ame khīcḍī jamīe chīe} — \textquotedblleft{}At night we eat khichdi.\textquotedblright{}
\end{quote}

\begin{cousinweb}
\textbf{\gu{ખીચડી}} (\emph{khīcḍī}) is a tadbhava. It comes from Sanskrit \emph{khiccā}, \textquotedblleft{}a dish of rice and peas\textquotedblright{}, extended with \emph{-ḍ-}, and CDIAL glosses the Gujarati word \textquotedblleft{}rice and lentils\textquotedblright{} (CDIAL 3880).
\end{cousinweb}

\begin{grammarlens}[title={What you've built}]
One more word for everyday life at A2.
\end{grammarlens}

\subsection*{Your turn}
\begin{itemize}
\raggedright
  \item \emph{Say it:} \emph{khīcḍī}
  \item \emph{Say it:} \emph{khīcḍī}, once more
  \item \emph{Say it:} \emph{rātre ame khīcḍī jamīe chīe} — At night we eat khichdi.
  \item \emph{Recall:} say the Gujarati for to roast, then the Gujarati for bhakhri, then say \emph{khīcḍī} again
\end{itemize}

\begin{tcolorbox}[breakable,colback=teal!4,colframe=teal!35!black,title={Before you move on}]
What does \gu{ખીચડી} mean? (\textquotedblleft{}Khichdi, rice and lentils\textquotedblright{}.) Say it once more.
\end{tcolorbox}
\section[\texorpdfstring{\gu{પાપડ} (pāpaḍ)}{પાપડ (pāpaḍ)}]{\gu{પાપડ} (pāpaḍ) — papad, poppadom}

\label{lesson:GU-C169-papad}



\begin{quote}
Before the new one: say the Gujarati for bhakhri, then the Gujarati for khichdi.
\end{quote}

\subsection*{You'll want to know: \gu{પાપડ}}
\textbf{\gu{પાપડ}} — \emph{pāpaḍ} — \textquotedblleft{}papad, poppadom\textquotedblright{}.

\begin{quote}
\textbf{\gu{પાપડ} \gu{પાતળો} \gu{છે}.} — \emph{pāpaḍ pātḷo che} — \textquotedblleft{}The papad is thin.\textquotedblright{}
\end{quote}

\begin{cousinweb}
\textbf{\gu{પાપડ}} (\emph{pāpaḍ}) is a tadbhava. It comes from Sanskrit \emph{parpaṭa}, \textquotedblleft{}a thin cake of rice or pease-meal\textquotedblright{}, and CDIAL lists the Gujarati \emph{pāpaḍ}, \textquotedblleft{}a thin cake\textquotedblright{} (CDIAL 7934).
\end{cousinweb}

\begin{grammarlens}[title={What you've built}]
One more word for everyday life at A2.
\end{grammarlens}

\subsection*{Your turn}
\begin{itemize}
\raggedright
  \item \emph{Say it:} \emph{pāpaḍ}
  \item \emph{Say it:} \emph{pāpaḍ}, once more
  \item \emph{Say it:} \emph{pāpaḍ pātḷo che} — The papad is thin.
  \item \emph{Recall:} say the Gujarati for bhakhri, then the Gujarati for khichdi, then say \emph{pāpaḍ} again
\end{itemize}

\begin{tcolorbox}[breakable,colback=teal!4,colframe=teal!35!black,title={Before you move on}]
What does \gu{પાપડ} mean? (\textquotedblleft{}Papad, poppadom\textquotedblright{}.) Say it once more.
\end{tcolorbox}
\section[\texorpdfstring{\gu{કડાઈ} (kaḍāī)}{કડાઈ (kaḍāī)}]{\gu{કડાઈ} (kaḍāī) — wok, deep pan}

\label{lesson:GU-C169-kadai}



\begin{quote}
Before the new one: say the Gujarati for khichdi, then the Gujarati for papad.
\end{quote}

\subsection*{You'll want to know: \gu{કડાઈ}}
\textbf{\gu{કડાઈ}} — \emph{kaḍāī} — \textquotedblleft{}wok, deep pan\textquotedblright{}.

\begin{quote}
\textbf{\gu{કડાઈમાં} \gu{તેલ} \gu{ગરમ} \gu{છે}.} — \emph{kaḍāīmā̃ tel garam che} — \textquotedblleft{}The oil in the wok is hot.\textquotedblright{}
\end{quote}

\begin{cousinweb}
\textbf{\gu{કડાઈ}} (\emph{kaḍāī}) is a tadbhava. It comes from Sanskrit \emph{kaṭāha}, \textquotedblleft{}boiler, saucepan with handles\textquotedblright{}, and CDIAL glosses the Gujarati word \textquotedblleft{}frying-pan\textquotedblright{} (CDIAL 2638).
\end{cousinweb}

\begin{grammarlens}[title={What you've built}]
One more word for everyday life at A2.
\end{grammarlens}

\subsection*{Your turn}
\begin{itemize}
\raggedright
  \item \emph{Say it:} \emph{kaḍāī}
  \item \emph{Say it:} \emph{kaḍāī}, once more
  \item \emph{Say it:} \emph{kaḍāīmā̃ tel garam che} — The oil in the wok is hot.
  \item \emph{Recall:} say the Gujarati for khichdi, then the Gujarati for papad, then say \emph{kaḍāī} again
\end{itemize}

\begin{tcolorbox}[breakable,colback=teal!4,colframe=teal!35!black,title={Before you move on}]
What does \gu{કડાઈ} mean? (\textquotedblleft{}Wok, deep pan\textquotedblright{}.) Say it once more.
\end{tcolorbox}
\section[\texorpdfstring{\gu{ઢાંકણું} (ḍhā̃kṇũ)}{ઢાંકણું (ḍhā̃kṇũ)}]{\gu{ઢાંકણું} (ḍhā̃kṇũ) — lid}

\label{lesson:GU-C169-dhaknu}



\begin{quote}
Before the new one: say the Gujarati for papad, then the Gujarati for wok.
\end{quote}

\subsection*{You'll want to know: \gu{ઢાંકણું}}
\textbf{\gu{ઢાંકણું}} — \emph{ḍhā̃kṇũ} — \textquotedblleft{}lid\textquotedblright{}.

\begin{quote}
\textbf{\gu{તપેલીનું} \gu{ઢાંકણું} \gu{ક્યાં} \gu{છે}?} — \emph{tapelīnũ ḍhā̃kṇũ kyā̃ che?} — \textquotedblleft{}Where is the lid of the pot?\textquotedblright{}
\end{quote}

\begin{cousinweb}
\textbf{\gu{ઢાંકણું}} (\emph{ḍhā̃kṇũ}) is a tadbhava, kin to \emph{ḍhā̃kvũ}, \textquotedblleft{}to cover\textquotedblright{}. CDIAL lists it, glossed \textquotedblleft{}cover, lid\textquotedblright{}, under \emph{*ḍhakk}, \emph{*ḍhaṅk}, \textquotedblleft{}cover\textquotedblright{} (CDIAL 5574).
\end{cousinweb}

\begin{grammarlens}[title={What you've built}]
One more word for everyday life at A2.
\end{grammarlens}

\subsection*{Your turn}
\begin{itemize}
\raggedright
  \item \emph{Say it:} \emph{ḍhā̃kṇũ}
  \item \emph{Say it:} \emph{ḍhā̃kṇũ}, once more
  \item \emph{Say it:} \emph{tapelīnũ ḍhā̃kṇũ kyā̃ che?} — Where is the lid of the pot?
  \item \emph{Recall:} say the Gujarati for papad, then the Gujarati for wok, then say \emph{ḍhā̃kṇũ} again
\end{itemize}

\begin{tcolorbox}[breakable,colback=teal!4,colframe=teal!35!black,title={Before you move on}]
What does \gu{ઢાંકણું} mean? (\textquotedblleft{}Lid\textquotedblright{}.) Say it once more.
\end{tcolorbox}
